\section{Sinusoidal error regression}
In this section, we test the effectiveness of \sindist{} in a simple linear regression setting by using it as a flexible distribution for the error terms. The linear model will hence be:
\[y_i = \beta_0 + \beta_1 x_i + \epsilon_i\]
-- where $\epsilon_i \iidsim \Sinu(a,d,s,k)$ and all model parameters $\beta_0, \beta_1, a, d, s, k$ are unknown -- in modelling various datasets. Our methodology will closely follow \cite{dean2009versatile}.

We will resort to numerical ML Estimation of model parameters. Observe that for SER, the data loglikelihood is:
\begin{align*}
	\ln L(\beta_0, \beta_1, a,d,s,k) &\equiv \ln \prod_{i=1}^n f(e_i; \beta_0, \beta_1, a,d,s,k)\\
	&= \sum_{i=1}^n \ln \psi\left(\frac{(y_i - \beta_0 - \beta_1 x_i) -a}{d}; s,k\right)
\end{align*}
which gives the MLE of the parameters as $\arg \max_{\beta_0, \beta_1, a,d,s,k} \sum_{i=1}^n \ln \psi\left(\frac{(y_i - \beta_0 - \beta_1 x_i) -a}{d}; s,k\right)$. Since analytical solutions are intractable, we will resort to a numerical optimization (L-BFGS-B) routine to extract ML estimates.

\subsection{Approximating usual Gauss-Markov setting}
First, we analyze the performance of \sindist{} in approximating a Gauss-Markov model with standard Normal error. Recollect in this regard that we already know that \sindist{} provides an arbitrarily well approximation of the Normal distribution for parameters in their limiting form. We repeated 50 iterations of the following experiment, with the arbitrary seed \texttt{set.seed(1234)} for reproducibility, for each of the sample sizes $n = 50, 100, 150, 200$ and collected performance metrics:

\begin{enumerate}
	\item Fix a sample size, say $n=50$.
	\item Generate 2 arbitrary coefficients $\beta_0 \sim U(10,20)$ and $\beta_1 \sim U(5,10)$, and a data column $x_i \sim U(0,1), \, i=1(1)n$. These stay fixed across all iterations.
	\item For every new iteration:
	\begin{enumerate}
		\item Simulate observational errors $\epsilon_i \iidsim \N(0,1), \, i=1(1)n$. Simulate the response as $y_i=\beta_0 + \beta_1*x_i + \epsilon$.
		\item Conduct Normal error regression, which is simply OLS regression. Note the coefficients.
		\item Conduct SER via numerical loglikelihood maximization. Note the coefficients.
	\end{enumerate}
	\item After all iterations are over, compare the NER coefficients against the SER coefficients.
	\item Move to the next sample size.
\end{enumerate}

Our findings are described in \autoref{tab:4NERvsSER}. Note that we obtained true $\beta_0 = 11.137$ and $\beta_1 = 8.111$ respectively.

\begin{table}[h!]
	\centering
	\begin{tabularx}{\linewidth}{X|X|X|X|X}
		Sample\newline Size  &  OLS $\hat\beta_0$ \newline Mean (Variance)  &  $SER \hat\beta_0$ \newline Mean (Variance)  &  OLS $\hat\beta_1$\newline Mean (Variance) &  SER $\hat\beta_1$\newline Mean (Variance)\\
		\hline
		10  & 11.175 (1.006) & 11.163 (1.024) & 7.970 (3.000) & 8.031 (3.208) \\
		25  & 11.103 (0.234) & 11.103 (0.234) & 8.219 (0.694) & 8.208 (0.690) \\
		50  & 11.363 (1.227) & 11.210 (0.110) & 8.008 (0.619) & 7.960 (00.549) \\
		100  & 11.155 (0.034) & 11.153 (.034) & 8.053 (0.144) & 8.069 (0.139) \\
		150  & 11.189 (0.019) & 11.191 (0.018) & 8.013 (0.067) & 8.009 (0.066) \\
		\hline
	\end{tabularx}
	\caption{Performance of SER against NER when errors were actually normally distributed}
	\label{tab:4NERvsSER}
\end{table}

We observe that due to severe computational complexities as described in \autoref{subsec:reg-disc}, there is no significant additional gain by implementing SER in a Gauss-Markov model. In fact, the estimate variances for $a,d,s,k$ are also unusably large for any practical or theoretical appreciation.

\begin{comment}
\subsection{Performance against actually Sinusoidal errors}
Next we test the applicability of the SER model when errors are truly distributed Sinusoidally. We conduct the following experiment 50 times for each of the sample sizes $n = 50, 250, 500, 750, 1000$ and collected performance metrics:

\begin{enumerate}
	\item Generate and fix arbitrary model parameters $\beta_0, \beta_1 \iidsim U(0,1)$ and $s, k \iidsim \U(0,10)$.
	\item For every $i=1(1)n$, generate and fix predictor points $x_i \iidsim \U(0,1)$,  and observational errors $\epsilon_i \iidsim \Sinu_0(0,1,s,k)$.
	\item Conduct NER as usual, and also SER via the loglikelihood maximization strategy described above.
\end{enumerate}
\end{comment}

\subsection{Discussion and Extensions}\label{subsec:reg-disc}
The optimization routine for likelihood estimation in a regression setup with a bounded support error-distribution such as the \sindist{} is particularly difficult. It contains, beyond the usual complexity of constrained maximization, the additional complexity of ensuring support coverage of the residuals at all times. i.e., for every $\beta_0, \beta_1$ chosen by the optimizer, a new set of estimates for the residuals are obtained, and consequently the optimizer's choice of $a,d$ must ensure that all errors lie within the support $[a, a+d]$. Hence the problem is manifold: where the region of constrained optimization itself varies with the optimizer's trajectory.

As a result of this, more complicated models such as SER on actually Sinusoidally distributed errors, and on arbitrary datasets with no distributional specification, are harder to implement. Safer, heuristic approaches such as modelling OLS residuals and choosing large enough supports for the error distribution, may however be tried.